\documentclass[11pt,a4paper]{article}

% Compile with: latexmk -pdf CV_Tinhinane_Amina_Azzouz.tex
% Academic facts checked against the public BIMSA profile, the previous CV,
% institutional course and seminar pages, the ICBS programme, and publisher
% and arXiv records. Updated 30 October 2026.
\usepackage[T1]{fontenc}
\usepackage[utf8]{inputenc}
\usepackage{lmodern}
\usepackage{microtype}
\usepackage[a4paper,left=21mm,right=21mm,top=19mm,bottom=19mm]{geometry}
\usepackage{xcolor}
\usepackage{array,tabularx}
\usepackage{enumitem}
\usepackage{titlesec}
\usepackage{fancyhdr}
\usepackage{lastpage}
\usepackage{hyperref}

\definecolor{accent}{HTML}{233D46}
\definecolor{muted}{HTML}{53616A}
\definecolor{rulecolor}{HTML}{CED6D8}
\hypersetup{colorlinks=true,urlcolor=accent,linkcolor=accent,
  pdftitle={Curriculum Vitae - Tinhinane Amina Azzouz},
  pdfauthor={Tinhinane Amina Azzouz},
  pdfsubject={Academic curriculum vitae, updated 30 October 2026},
  pdfkeywords={p-adic differential equations, non-Archimedean geometry, spectral theory}}
\urlstyle{same}
\setlength{\parindent}{0pt}
\setlength{\parskip}{2pt}
\setlength{\tabcolsep}{0pt}
\renewcommand{\arraystretch}{1.0}
\AtBeginDocument{\fontsize{10.5}{12.8}\selectfont}
\raggedbottom
\setlist[itemize]{leftmargin=1.15em,topsep=2pt,itemsep=3.5pt,parsep=0pt}
\setlist[enumerate]{leftmargin=1.7em,topsep=2pt,itemsep=7pt,parsep=0pt}
\titleformat{\section}{\large\bfseries\color{accent}}{}{0pt}{}[\vspace{-3pt}\color{rulecolor}\titlerule]
\titlespacing*{\section}{0pt}{11pt}{5pt}
\pagestyle{fancy}
\fancyhf{}
\renewcommand{\headrulewidth}{0pt}
\renewcommand{\footrulewidth}{0pt}
\fancyfoot[L]{\small\color{muted}Tinhinane Amina Azzouz\enspace\textperiodcentered\enspace Curriculum vitae}
\fancyfoot[R]{\small\color{muted}\thepage\,/\,\pageref*{LastPage}}
\setlength{\footskip}{10mm}

\newcommand{\entry}[3]{%
  \begin{tabularx}{\textwidth}{@{}>{\raggedright\arraybackslash}p{32mm}@{\hspace{4mm}}>{\raggedright\arraybackslash}X@{}}%
  \textcolor{muted}{#1} & \textbf{#2}\newline #3
  \end{tabularx}\par\vspace{3pt}}
\newcommand{\shortentry}[2]{%
  \begin{tabularx}{\textwidth}{@{}>{\raggedright\arraybackslash}p{32mm}@{\hspace{4mm}}>{\raggedright\arraybackslash}X@{}}\textcolor{muted}{#1} & #2\end{tabularx}\par\vspace{3pt}}
\newcommand{\arxiv}[1]{\href{https://arxiv.org/abs/#1}{arXiv:#1}}
\newcommand{\doi}[1]{\href{https://doi.org/#1}{doi:#1}}

\begin{document}

{\fontsize{26}{29}\selectfont\bfseries\color{accent}Tinhinane Amina Azzouz}\par
\vspace{3pt}
{\large Curriculum vitae}\hfill{\small\color{muted}Updated 30 October 2026}\par
\vspace{7pt}
\textbf{Assistant Professor of Mathematics}\par
Beijing Institute of Mathematical Sciences and Applications (BIMSA)\par
Huairou District, Beijing 101408, China\quad\textperiodcentered\quad Office A12-202\par
\href{mailto:azzouzta@bimsa.cn}{azzouzta@bimsa.cn}\quad\textperiodcentered\quad
\href{https://www.bimsa.cn/detail/tinhinaneazzouz.html}{BIMSA profile}\quad\textperiodcentered\quad
\href{https://tinaamina.github.io/}{Personal webpage}\par

\section{Research interests}
$p$-adic differential equations; non-Archimedean spectral theory; ultrametric functional analysis; Berkovich analytic geometry.

\section{Academic appointments}
\entry{Jan.\ 2024--present}{Assistant Professor}{Beijing Institute of Mathematical Sciences and Applications, Beijing, China.}
\entry{2021--2023}{Postdoctoral Researcher}{Yau Mathematical Sciences Center, Tsinghua University / BIMSA, Beijing, China.}
\entry{Nov.\ 2018--Dec.\ 2020}{Assistant Professor}{Universit\'e Alger 1, Algiers, Algeria.}
\entry{Sept.--Nov.\ 2018}{Postdoctoral Researcher}{Max Planck Institute for Mathematics, Bonn, Germany.}
\entry{Oct.\ 2017--Aug.\ 2018}{Temporary Teaching and Research Assistant (ATER)}{Universit\'e Grenoble Alpes, Grenoble, France.}

\section{Education}
\entry{2014--2018}{Ph.D. in Mathematics}{Universit\'e de Montpellier, Institut Montpelli\'erain Alexander Grothendieck, France.\newline
Thesis: \emph{Spectrum of $p$-adic differential equations}.\newline
Supervisors: Andrea Pulita and J\'er\^ome Poineau. Defended 11 June 2018.}
\entry{2012--2014}{M.Sc. in Mathematics and Applications}{Fundamental Mathematics, Universit\'e Pierre et Marie Curie, Paris, France. With honours.}
\entry{2008--2011}{B.Sc. in Mathematics and Computer Science}{Algebra and Cryptography, Universit\'e des Sciences et de la Technologie Houari Boumediene (USTHB), Algiers, Algeria. With honours.}

\section{Funding and awards}
\shortentry{2023}{\textbf{Research Fund for International Young Scientists}, National Natural Science Foundation of China; YMSC, Tsinghua University.}
\shortentry{2014--2017}{\textbf{Funded doctoral position with teaching}, Universit\'e de Montpellier.}
\shortentry{2012--2014}{\textbf{Paris Graduate School of Mathematics international programme}, laureate; Fondation Sciences Math\'ematiques de Paris.}

\section{Academic service}
\shortentry{Nov.\ 2019--June 2020}{\textbf{Deputy Director in charge of Teaching}, Department of Mathematics and Computer Science, Universit\'e Alger 1.}

\newpage
\section{Publications}
\begin{enumerate}
\item Tinhinane Amina Azzouz. \emph{Spectrum of $p$-adic linear differential equations II: Variation of the spectrum}. \textbf{Journal of the London Mathematical Society} \textbf{114}(4) (2026), e70718.\par
\doi{10.1112/jlms.70718}. \arxiv{2303.06014}.
\item Tinhinane Amina Azzouz. \emph{Spectrum of $p$-adic linear differential equations I: The shape of the spectrum}. \textbf{Selecta Mathematica} \textbf{30}, article 13 (2024).\par
\doi{10.1007/s00029-023-00904-4}. \arxiv{2111.03548}.
\item Tinhinane Amina Azzouz. \emph{Spectrum of a linear differential equation over a field of formal power series}. \textbf{Journal of Number Theory} \textbf{231} (2022), 139--157.\par
\doi{10.1016/j.jnt.2020.11.021}. \arxiv{1808.02382}.
\item Tinhinane Amina Azzouz. \emph{Spectrum of a linear differential equation with constant coefficients}. \textbf{Mathematische Zeitschrift} \textbf{296} (2020), 1613--1644.\par
\doi{10.1007/s00209-020-02482-z}. \arxiv{1802.07306}.
\end{enumerate}

\section{Preprints}
Tinhinane Amina Azzouz, Mainak Ghosh and Sebastien Palcoux.
\emph{Local Ocneanu rigidity for separable algebra objects} (2026).
\arxiv{2609.10440}.

\section{Work in progress}
\begin{itemize}
\item Wild finite \'etale coverings of $p$-adic annuli, with Andrea Pulita.
\item Spectral methods and decompositions for $p$-adic differential modules.
\item Techniques for computing the spectrum of an ultrametric operator.
\end{itemize}

\section{Teaching}
\textbf{Courses at BIMSA}\par\vspace{2pt}
\shortentry{Oct.\ 2026--Jan.\ 2027}{\href{https://www.bimsa.cn/research_detail/RamTheofLocFie.html}{\textbf{Ramification Theory of Local Fields}}.}
\shortentry{Oct.--Dec.\ 2025}{\href{https://www.bimsa.cn/research_detail/InttoBerspa.html}{\textbf{Introduction to Berkovich spaces}}.}
\shortentry{Feb.--May 2025}{\href{https://www.bimsa.cn/research_detail/Inttop-adifequIII.html}{\textbf{Introduction to $p$-adic differential equations III}}.}
\shortentry{Oct.\ 2024--Jan.\ 2025}{\href{https://www.bimsa.cn/research_detail/Inttop-adifequII.html}{\textbf{Introduction to $p$-adic differential equations II}}.}
\shortentry{Mar.--June 2024}{\href{https://www.bimsa.cn/research_detail/Inttop-adifequ.html}{\textbf{Introduction to $p$-adic differential equations}}.}
\vspace{5pt}
\textbf{Earlier teaching}\par\vspace{2pt}
\shortentry{Aug.\ 2022}{\textbf{Introduction to $p$-adic differential equations}, Sanya Mathematics Summer School, Tsinghua Sanya International Mathematics Forum, Sanya, China.}
\shortentry{Oct.\ 2015--Aug.\ 2016}{Teaching assignment, Universit\'e Grenoble Alpes, France.}
\shortentry{Oct.\ 2014--Aug.\ 2015}{Teaching assignment, Universit\'e de Montpellier, France.}

\newpage
\section{Selected talks}
\shortentry{14 Aug.\ 2026}{\emph{The spectrum of $p$-adic linear differential equations}. Invited lecture, International Congress of Basic Science (ICBS), Beijing, China.}
\shortentry{6 \& 13 Nov.\ 2025}{\emph{The variation of the radii of convergence of linear $p$-adic differential equation}. Number Theory Lunch Seminar, BIMSA, Beijing.}
\shortentry{27 May 2024}{\emph{$p$-adic linear differential equations}. BIMSA Member Seminar, Beijing.}
\shortentry{27 Sept.\ 2023}{\emph{$p$-adic differential equations and radii of convergence}. Number Theory Lunch Seminar, BIMSA, Beijing.}
\shortentry{4 July 2023}{\emph{Spectrum of $p$-adic linear differential equations}. Tenth Congress of Romanian Mathematicians, Romania.}
\shortentry{18 May 2023}{\emph{The spectral theory in the sense of Berkovich}. Number Theory Lunch Seminar, BIMSA, Beijing.}
\shortentry{28 Oct.\ 2022}{\emph{Finite morphism and radii of convergence of a $p$-adic differential equation}. Number Theory Lunch Seminar, BIMSA, Beijing.}
\shortentry{June 2022}{\emph{Spectrum of $p$-adic differential equations}. BIMSA--YMSC Tsinghua Number Theory Seminar, Beijing.}
\shortentry{July 2017}{\emph{Spectral theory and ultrametric differential equations}. Students' Conference on Tropical and Non-Archimedean Geometry, University of Regensburg, Germany.}
\shortentry{June 2016}{\emph{Spectre d'un syst\`eme diff\'erentiel aux coefficients constants}. \'Ecole jeunes chercheurs en th\'eorie des nombres, Clermont-Ferrand, France.}
\shortentry{2014--2018}{Five talks in Ph.D. students' seminars.}

\section{Conference and school participation}
\shortentry{July 2019}{\emph{Perfectoids}, Centre Henri Lebesgue, Rennes, France.}
\shortentry{July 2018}{\emph{Berkovich Spaces 30 Years}, Institut Henri Poincar\'e, Paris, France.}
\shortentry{July 2017}{\emph{Berkovich Spaces, Tropical Geometry and Model Theory}, Universidad de los Andes, Bogot\'a, Colombia.}
\shortentry{Mar.\ 2017}{\emph{$p$-adic Analytic Geometry and Differential Equations}, CIRM, Marseille, France.}
\shortentry{Aug.\ 2015}{Students' Conference on Tropical and Non-Archimedean Geometry, University of Regensburg, Germany.}
\shortentry{June 2015}{\emph{G\'eom\'etrie et arithm\'etique sur les corps locaux et globaux}, Universit\'e de Caen, France.}
\shortentry{June 2015}{Conference in honour of Alexandre Grothendieck, Universit\'e de Montpellier, France.}
\shortentry{June 2014}{\emph{Special Week: Berkovich Spaces and Applications}, Universit\'e de Strasbourg, France.}

\section{Languages}
Arabic (native); French (fluent); English (fluent).

\end{document}
